\documentclass[11pt,a4paper]{article}
\usepackage[utf8]{inputenc}
\usepackage[vietnamese]{babel}
\usepackage{amsmath,amssymb}
\usepackage{geometry}
\geometry{a4paper, total={180mm,260mm}, left=15mm, right=15mm, top=18mm, bottom=18mm}
\usepackage{tikz}
\usetikzlibrary{shapes.geometric, arrows.meta, positioning, calc}
\usepackage{tcolorbox}
\tcbuselibrary{skins,breakable}
\usepackage{xcolor}
\usepackage{multicol}
\usepackage{array}
\usepackage{fontawesome}
\usepackage{fancyhdr}

\pagestyle{fancy}
\fancyhf{}
\renewcommand{\headrulewidth}{0.5pt}
\renewcommand{\footrulewidth}{0.5pt}
\fancyhead[L]{\small\bfseries\color{teal}\faBook\ SỔ TAY KIẾN THỨC TOÁN 3}
\fancyhead[R]{\small\bfseries\color{teal}Lớp Toán Thầy Lĩnh -- loptoanthaylinh.vn}
\fancyfoot[C]{\small\color{gray}-- Sổ tay kiến thức: Bài 80. Ôn tập bảng số liệu và khả năng xảy ra --}

\definecolor{primary}{HTML}{005F73}
\definecolor{secondary}{HTML}{0A9396}
\definecolor{accent}{HTML}{EE9B00}
\definecolor{danger}{HTML}{AE2012}
\definecolor{boxbg}{HTML}{F4F6F8}

\begin{document}

% --- BANNER TIÊU ĐỀ ---
\begin{tcolorbox}[
  enhanced,
  colback=white,
  colframe=primary,
  arc=3mm,
  boxrule=1.2pt,
  left=4mm, right=4mm, top=2.5mm, bottom=2.5mm,
  drop shadow=black!10!white
]
  \centering
  {\Large\bfseries\color{primary}\faBookmark\ BÀI 80: ÔN TẬP BẢNG SỐ LIỆU VÀ KHẢ NĂNG XẢY RA}\\[1.5mm]
  {\small\itshape\color{black!75}Ôn tập cuối năm: Kĩ năng phân tích bảng số liệu thống kê và 3 mức độ xác suất Chắc chắn - Có thể - Không thể trên 1 trang A4}
\end{tcolorbox}

\vspace{-2mm}


\noindent
\begin{minipage}[t]{0.49\textwidth}\vspace{0pt}
\begin{tcolorbox}[
  enhanced,
  colback=blue!4!white,
  colframe=blue!75!black,
  arc=2mm, boxrule=0.8pt,
  left=2mm, right=2mm, top=1.8mm, bottom=1.8mm,
  title={\footnotesize\bfseries\color{white}\faKey\ 1. Đọc và phân tích Bảng số liệu}
]
{\footnotesize
\begin{itemize}\setlength{\itemsep}{1.5pt}\setlength{\leftskip}{-2.5mm}
  \item Đọc kĩ tiêu đề hàng và tiêu đề cột của bảng số liệu.
  \item Xác định số liệu tương ứng với từng đối tượng.
  \item \textbf{Các dạng câu hỏi phân tích thường gặp:}\\
  1. Đối tượng nào có số lượng nhiều nhất? Ít nhất?\\
  2. Hai đối tượng hơn kém nhau bao nhiêu đơn vị?\\
  3. Cả tập hợp có tất cả bao nhiêu đối tượng (tính tổng)?
\end{itemize}
}
\end{tcolorbox}
\end{minipage}
\hfill
\begin{minipage}[t]{0.49\textwidth}\vspace{0pt}
\begin{tcolorbox}[
  enhanced,
  colback=teal!4!white,
  colframe=teal!75!black,
  arc=2mm, boxrule=0.8pt,
  left=2mm, right=2mm, top=1.8mm, bottom=1.8mm,
  title={\footnotesize\bfseries\color{white}\faSliders\ 2. Ba mức độ khả năng của sự kiện}
]
{\footnotesize
\begin{itemize}\setlength{\itemsep}{1.5pt}\setlength{\leftskip}{-2.5mm}
  \item \textbf{Chắc chắn:} Luôn luôn xảy ra trong mọi trường hợp (tất cả các đối tượng đều thỏa mãn).
  \item \textbf{Có thể:} Có trường hợp xảy ra, có trường hợp không xảy ra (chỉ một số đối tượng thỏa mãn).
  \item \textbf{Không thể:} Không bao giờ xảy ra trong bất kì trường hợp nào (không có đối tượng nào thỏa mãn).
\end{itemize}
}
\end{tcolorbox}
\end{minipage}

\vspace{-2.5mm}
% --- KHỐI 3: SƠ ĐỒ TƯ DUY TIKZ ---
\noindent
\begin{tcolorbox}[
  enhanced,
  colback=orange!4!white,
  colframe=orange!80!black,
  arc=2mm, boxrule=0.8pt,
  left=2mm, right=2mm, top=1.5mm, bottom=1.5mm,
  title={\footnotesize\bfseries\color{white}\faSitemap\ SƠ ĐỒ TƯ DUY: HỆ THỐNG THỐNG KÊ VÀ XÁC SUẤT LỚP 3}
]
\centering
\begin{tikzpicture}[
  box/.style={draw, rounded corners=2pt, align=center, font=\scriptsize\bfseries, inner sep=2.5pt, minimum height=0.75cm},
  arr/.style={-{Stealth[length=2mm]}, thick}
]
  \node[box, fill=purple!15, draw=purple!80!black, text width=2.6cm] (t1) at (0, 0) {BẢNG SỐ LIỆU\\Đọc theo hàng \& cột\\Phân tích dữ liệu};
  \node[box, fill=green!20, draw=green!70!black, text width=2.6cm] (t2) at (3.5, 0) {CHẮC CHẮN\\$100\%$ xảy ra\\Mọi trường hợp};
  \node[box, fill=blue!15, draw=blue!80!black, text width=2.6cm] (t3) at (7.0, 0) {CÓ THỂ\\Có thể xảy ra / không\\Yếu tố ngẫu nhiên};
  \node[box, fill=red!15, draw=red!70!black, text width=2.6cm] (t4) at (10.5, 0) {KHÔNG THỂ\\$0\%$ xảy ra\\Hoàn toàn không có};

  \draw[arr, purple!80!black] (t1) -- (t2);
  \draw[arr, teal!80!black] (t2) -- (t3);
  \draw[arr, orange!80!black] (t3) -- (t4);
\end{tikzpicture}
\end{tcolorbox}

\vspace{-2.5mm}
% --- KHỐI 4: VÍ DỤ & MẸO TRÁNH BẪY ---
\noindent
\begin{minipage}[t]{0.49\textwidth}\vspace{0pt}
\begin{tcolorbox}[
  enhanced,
  colback=green!4!white,
  colframe=green!60!black,
  arc=2mm, boxrule=0.8pt,
  left=2mm, right=2mm, top=1.8mm, bottom=1.8mm,
  title={\footnotesize\bfseries\color{white}\faCheckCircle\ Ví dụ thực hành chuẩn mực}
]
{\footnotesize
\textbf{Trong hộp có 4 bút mực xanh và 1 bút mực tím:}\\
-- Rút được bút mực xanh: \textbf{Có thể}.\\
-- Rút được bút mực đen: \textbf{Không thể}.\\
-- Rút được chiếc bút có màu: \textbf{Chắc chắn}.\\
\textbf{Bảng điểm kiểm tra học kì:}\\
Số bạn đạt điểm $10$ là $12$ bạn, điểm $9$ là $15$ bạn. Số bạn đạt điểm $9$ và $10$ là: $12 + 15 = \mathbf{27\text{ bạn}}$.
}
\end{tcolorbox}
\end{minipage}
\hfill
\begin{minipage}[t]{0.49\textwidth}\vspace{0pt}
\begin{tcolorbox}[
  enhanced,
  colback=red!4!white,
  colframe=red!70!black,
  arc=2mm, boxrule=0.8pt,
  left=2mm, right=2mm, top=1.8mm, bottom=1.8mm,
  title={\footnotesize\bfseries\color{white}\faExclamationTriangle\ Mẹo tránh bẫy trắc nghiệm}
]
{\footnotesize
\begin{itemize}\setlength{\itemsep}{1.5pt}\setlength{\leftskip}{-2.5mm}
  \item \textbf{Bẫy vội vàng chọn "Chắc chắn":}\\
  Dù số lượng bút xanh rất nhiều ($4$ chiếc) và bút tím chỉ có $1$ chiếc, việc rút được bút xanh vẫn chỉ là \textbf{"Có thể"}, không bao giờ được chọn "Chắc chắn"!
  \item Đọc kĩ danh sách đối tượng có trong đề bài.
\end{itemize}
}
\end{tcolorbox}
\end{minipage}


\end{document}
